\section{Conclusions}
\label{sec:conclusions}

\begin{enumerate}
\item
An analytical model has been developed for branching initiation at the
corner of a V-shaped interfacial shear crack. The singular field of the
prescribed stationary crack is used as the outer field for a small tensile
process zone oriented along the bisector of one of the materials.
Closed-form relations are obtained for the equilibrium zone length $d_i$,
opening $\delta_i$, work $W_i$ done against the resistance to separation,
and the zone-related energy parameter $J_i=\dd W_i/\dd d_i$.

\item
Physical admissibility of the local process zone requires both compressive
contact of the V-shaped interfacial crack faces, $Cg_2<0$, and a tensile
normal field in the material containing the zone, $CQ_i>0$. For the
parameters examined, these conditions can be satisfied simultaneously in
only one of the two materials for a given angle and loading sign. The
interface angle and elastic contrast shift the admissible ranges and the
equilibrium zone characteristics.

\item
The presence of the process zone weakens the corner singularity of the
configuration but does not eliminate it. On the physically admissible intervals for the baseline
material pair, $-1/2<\lambda_i<0$. The residual field is therefore weaker
than the classical $r^{-1/2}$ crack singularity, and its local
configurational elastic energy flux vanishes as the contour contracts to
the corner. Consequently, no separate finite elastic additive term is required for the
zone-related energy parameter $J_i$ in the local branching-initiation
criterion for this corner problem.

\item
For fixed geometry, the zone opening and the zone-related energy parameter
are uniquely related through the same equilibrium zone length:
$J_i=\chi_i\sigma_i\delta_i$. The deformation criterion
$\delta_i=\delta_{i\mathrm{c}}$ and the energy criterion
$J_i=J_{i\mathrm{c}}$ therefore identify the same critical state only
when the independently specified critical parameters are calibrated
consistently,
$J_{i\mathrm{c}}=\chi_i\sigma_i\delta_{i\mathrm{c}}$. Once either
critical material parameter is supplied independently, the model determines
the corresponding critical external load for branching initiation.

\item
Quantitative use of the model relies on the small-scale assumption
$d_i\ll l$; $d_i/l<0.1$ is adopted as the principal applicability range,
whereas portions of the curves reaching approximately $d_i/l=0.18$ should
be interpreted as extrapolations of the local approximation. Branching
from the corner and interfacial propagation from the outer tips of the
V-shaped crack are competing local mechanisms rather than prescribed
sequential stages. Because the formulation is local, it can be extended
to developed V-shaped interfacial cracks by replacing the outer-field
amplitude with the corresponding corner-field amplitude. A related
application is splitting of a solid by a sharp wedge.
\end{enumerate}
